\chapter{Plumbing}
\label{ch:plumbing}

\index{plumbing}
\index{plumber}

When you right-click (\button{3}) on text in acme, the \emph{plumber}
decides what to do with it.  Is it a filename?  Open it.  A URL?
Launch a browser.  A compiler error?  Jump to the source line.  The
plumber is a pattern-matching system that routes text to the appropriate
action.

This standalone acme has an embedded plumber---no external daemon or
configuration is needed.  It comes with sensible defaults and supports
custom rules.

\section{How Plumbing Works}

When you right-click on text, acme:

\begin{enumerate}
\item Identifies the text under the click (expands to word/filename
      boundaries).
\item Constructs a \emph{plumb message} containing the text, its type,
      the source directory, and any attributes (like an address).
\item Matches the message against plumbing rules in order.
\item The first matching rule determines the action: open in the editor,
      start an external program, or search.
\end{enumerate}

\section{Default Rules}
\label{sec:defaultrules}

\index{plumbing!default rules}

The built-in rules handle common patterns:

\begin{longtable}{@{}lp{3.2in}@{}}
\toprule
Pattern & Action \\
\midrule
\endhead
URLs (\texttt{http://}, \texttt{https://}, etc.) &
  Open in web browser. \\
Image files (\texttt{.jpg}, \texttt{.png}, \texttt{.gif}, etc.) &
  Open in image viewer. \\
Document files (\texttt{.pdf}, \texttt{.ps}, \texttt{.dvi}) &
  Open in document viewer. \\
\texttt{file:line.col,line.col} &
  Open file at line and column range. \\
\texttt{file:line.col} &
  Open file at line and column. \\
\texttt{file:line} or \texttt{file:/regexp/} &
  Open file at address. \\
Plain filenames &
  Open file in acme. \\
\texttt{<file.h>} style includes &
  Search \file{/usr/include} and \file{/usr/local/include}. \\
\bottomrule
\end{longtable}

\subsection{Platform Differences}
\index{plumbing!platform}

The command used to open URLs, images, and PDFs varies by platform:

\begin{itemize}
\item \textbf{Linux:} \cmd{xdg-open}
\item \textbf{macOS:} \cmd{open}
\item \textbf{Windows (MSYS2):} \cmd{cygstart}
\end{itemize}

\section{Custom Rules}
\label{sec:customrules}

\index{plumbing!custom rules}
\index{plumbing rules file@\file{\textasciitilde/lib/plumbing}}

To override or extend the default rules, create the file
\file{\textasciitilde/lib/plumbing}.  If this file exists, it
\emph{replaces} the built-in defaults entirely.

\subsection{Rule Syntax}

A plumbing rule is a sequence of patterns and actions, separated by
blank lines.  Each rule consists of:

\begin{lstlisting}[frame=none,basicstyle=\small\ttfamily]
# Optional comment
type is text
data matches '<regexp>'
arg isfile $1
plumb to <port>
plumb start <command> $file
\end{lstlisting}

\paragraph{Pattern lines} test properties of the plumb message:

\begin{longtable}{@{}lp{3.2in}@{}}
\toprule
Pattern & Meaning \\
\midrule
\endhead
\texttt{type is text} & The message type is ``text''. \\
\texttt{data matches '\textit{re}'} & The message data matches the regex.
  Capturing groups set \texttt{\$1}, \texttt{\$2}, etc. \\
\texttt{data set \$file} & Set the data to the matched file path. \\
\texttt{arg isfile \$\textit{n}} & Check that the captured group is an
  existing file path.  Sets \texttt{\$file} if so. \\
\texttt{attr add addr=\$\textit{n}} & Add an attribute (e.g., a line
  number address). \\
\bottomrule
\end{longtable}

\paragraph{Action lines} specify what to do:

\begin{longtable}{@{}lp{3.2in}@{}}
\toprule
Action & Meaning \\
\midrule
\endhead
\texttt{plumb to edit} & Open in acme editor. \\
\texttt{plumb to web} & Open as a URL. \\
\texttt{plumb to image} & Open as an image. \\
\texttt{plumb client \$editor} & Open in the acme editor (for ``edit''
  port). \\
\texttt{plumb start \textit{cmd} \$file} & Run an external command. \\
\bottomrule
\end{longtable}

\subsection{Variables}

\begin{longtable}{@{}lp{3.5in}@{}}
\toprule
Variable & Meaning \\
\midrule
\endhead
\texttt{\$0} & The entire matched data. \\
\texttt{\$1}, \texttt{\$2}, \ldots & Captured groups from \texttt{data matches}. \\
\texttt{\$file} & The resolved file path (set by \texttt{arg isfile}). \\
\texttt{\$editor} & The editor variable (set at top of rules). \\
\bottomrule
\end{longtable}

\subsection{Example Custom Rules}

\paragraph{Open Go import paths:}
\begin{lstlisting}[frame=none,basicstyle=\small\ttfamily]
# Go import paths -> go doc
type is text
data matches '[a-zA-Z0-9./\-]+'
data matches '([a-z]+\.[a-z]+/[a-zA-Z0-9./\-]+)'
plumb to web
plumb start xdg-open https://pkg.go.dev/$1
\end{lstlisting}

\paragraph{Open Python files with line numbers from tracebacks:}
\begin{lstlisting}[frame=none,basicstyle=\small\ttfamily]
type is text
data matches 'File "([^"]+)", line ([0-9]+)'
arg isfile $1
data set $file
attr add addr=$2
plumb to edit
plumb client $editor
\end{lstlisting}

\paragraph{Open man pages:}
\begin{lstlisting}[frame=none,basicstyle=\small\ttfamily]
type is text
data matches '([a-zA-Z0-9_\-]+)\(([0-9])\)'
plumb to none
plumb start xterm -e man $2 $1
\end{lstlisting}

\section{Debugging Plumbing}

If right-clicking does not produce the expected result:

\begin{enumerate}
\item Check \file{\textasciitilde/lib/plumbing} for syntax errors.
\item Remember that custom rules \emph{replace} the defaults---include
      the standard rules in your file if you want to keep them.
\item Rules are matched in order; the first match wins.
\item File existence checks (\texttt{arg isfile}) use the source
      window's directory for relative paths.
\end{enumerate}
